

Let $E_{12}\to E$ be the connected cyclic cover defined by
reduction modulo twelve of the positive meridian character.
The circle fibration of \ref{maps:src:E.2.13} may be chosen with fibre
\[
 Y=(E_{12}\times T^3)\ \cup_{T^8}\ (T^3\times E_{12}),
 \tag{E.S29}\label{maps:eq:E.S29}
\]
with the following exact boundary marking. Write
$\theta_i'$ for the lifted normal circle and retain the
two regular-torus markings. Then
\[
       \theta_1'=\gamma_2-\gamma_1,\qquad
       \theta_2'=\gamma_1-\gamma_2.
 \tag{E.S30}\label{maps:eq:E.S30}
\]
There is a return representative $\psi$ whose twelfth power
is the separating torus twist
\[
 \begin{gathered}
 T(x,s)=(x+\eta(s)(g_1+g_2),s)\quad\hbox{on }T^8\times[-1,1],
 \\
 T=\operatorname{id}\text{ elsewhere}.
 \end{gathered}
 \tag{E.S31}\label{maps:eq:E.S31}
\]
Here $\eta$ is zero near $-1$ and one near $1$, and circle
coordinates have period one. In particular \eqref{maps:eq:E.S31} is the
identity near both ends of its collar. Its rational
cohomological variation satisfies
\[
 (T^*-I)^2=0,\qquad
 \operatorname{rank}(T^*-I)|_{H^j(Y;\mathbb Q)}=
 \begin{cases}1&j=4,5,\\0&j\ne4,5.\end{cases}
 \tag{E.S32}\label{maps:eq:E.S32}
\]
The rational Betti numbers of the actual fibre are
\[
       (b_0,\ldots,b_9)(Y)
          =(1,12,121,407,693,693,407,121,12,1).
 \tag{E.S33}\label{maps:eq:E.S33}
\]
Every eigenvalue of $\psi^*$ is a twelfth root of unity.
At eigenvalue one there is exactly one Jordan block of
size two in each of degrees four and five; all its other
blocks, and all blocks at other eigenvalues, have size one.
Thus the infinite order is already visible on ordinary
rational cohomology. The smooth complement $E_{12}$ in
\eqref{maps:eq:E.S29} remains the actual finite cover of $E$, not a
replacement inferred from its group or homology.

\begin{proof}
\emph{Adapting the fibration to the actual splice.}
Cut the base along the circle used in \ref{maps:src:E.2.6}, with the two
critical triples on opposite sides. Over the first disc the
second factor is a regular marked $T^4$, and over the second
disc the first factor is regular. Let
$\xi_E:E\to\mathbb R/\mathbb Z$ represent the positive
meridian class. Its restriction to the boundary can be
chosen exactly
\[
                         \xi_E=\theta+12\gamma.
 \tag{E.S34}\label{maps:eq:E.S34}
\]
Indeed this is its integral restriction by \ref{maps:src:E.2.10}; homotopy
extension makes the representative equal to that linear
map on a collar. The two maps
\[
        \rho_A=\xi_{E_1}-12\gamma_2,
             \qquad \rho_B=12\gamma_1-\xi_{E_2}
 \tag{E.S35}\label{maps:eq:E.S35}
\]
agree exactly on the common boundary, where
$\theta_2=-\theta_1$. Smooth them with this common product
collar. The regular-factor circle derivative is $-12$ on
the first piece and $12$ on the second, so they give a
proper submersion throughout, including both cusp pieces.
These are also permissible choices in E.2.13: on the
quotient $N$, an invariant closed one-form differs from
these representatives only by the differential of a
single-valued function on each of the two pieces. Absorb
that function by translation in the unused regular circle.
Thus this construction uses the same actual quotient and
primitive character.

A level of $\xi_E-12\gamma_2$ is the pullback of the
degree-twelve circle cover along $\xi_E$, times the three
remaining regular circles. The pullback is $E_{12}$.
Its boundary is connected because the normal meridian has
character one. The covering of that boundary has
$\theta=12\theta'$ and unchanged $\gamma,u,w,\delta$.
Equation \eqref{maps:eq:E.S34} in the level condition gives precisely
\eqref{maps:eq:E.S30}. This proves \eqref{maps:eq:E.S29} with its actual smooth
boundary maps.

Choose the vector field integrating \eqref{maps:eq:E.S35} to be
$-g_2/12$ on the first piece off its collar and $g_1/12$
on the second. On the common regular collar interpolate by
\[
                     W_s=\frac{\eta(s)g_1-(1-\eta(s))g_2}{12}.
 \tag{E.S36}\label{maps:eq:E.S36}
\]
It has $d\rho(W_s)=1$, and it preserves the collar
coordinate $s$ and the original base point. Its flow
therefore defines a return $\psi$. At time twelve it is
the identity off the collar, since the two indicated
regular circles have period one. On the collar it
translates by $\eta(s)g_1-(1-\eta(s))g_2$; removing the
integral translation $-g_2$ gives \eqref{maps:eq:E.S31}. This comparison
uses a globally regular factor on each piece, so there
is no continuation through a singular cap to justify.

\emph{The variation map and the invariant boundary calculation.}
Put $D=T^8$ and let $r:H^j(Y;\mathbb Q)\to H^j(D;\mathbb Q)$
be restriction. Denote the Mayer--Vietoris connecting map
by $\partial:H^{j-1}(D;\mathbb Q)\to H^j(Y;\mathbb Q)$.
Writing $v=g_1+g_2$, the collar pullback of a translation
invariant form is
\[
       T^*\omega=\omega+d\eta\wedge\iota_v\omega.
\]
Average its restriction to the torus and change the global
representative by an exact form to make this equality
applicable on the collar. The difference represents the
connecting class, and hence, with the connecting sign
chosen by the increasing collar coordinate,
\[
                          T^*-I=\partial\,\iota_v\,r.
 \tag{E.S37}\label{maps:eq:E.S37}
\]
As $r\partial=0$, its square is zero.

No unknown part of $H^*(E_{12})$ is needed to compute
\eqref{maps:eq:E.S37}. Its cyclic deck action is a translation on the
boundary torus, hence acts trivially on boundary cohomology.
Every noninvariant deck eigenspace therefore restricts to
zero. Rational transfer identifies the invariant part
with $H^*(E)$, whose marked restriction is \ref{maps:src:E.2.6}. Its
pullback to \eqref{maps:eq:E.S30} is
\[
 R_{k,1}(a)=\phi_k(a)+12(\gamma_2-\gamma_1)a,
 \qquad
 R_{k,2}(b)=\phi_k(b)+12(\gamma_1-\gamma_2)b.
 \tag{E.S38}\label{maps:eq:E.S38}
\]
Retain the other factor's three covectors $u,w,\delta$,
and take exterior products to obtain the two subspaces
$U_j^k\subset\bigwedge^k\mathbb Q^8$.
The image of $r$ is exactly $U_1^k\cap U_2^k$, and the
kernel of $\partial$ is $U_1^{k-1}+U_2^{k-1}$.
Consequently the rank in \eqref{maps:eq:E.S32} is exactly the rank of
\[
 U_1^k\cap U_2^k\ \stackrel{\iota_v}{\longrightarrow}
 \bigwedge^{k-1}\mathbb Q^8/(U_1^{k-1}+U_2^{k-1}).
 \tag{E.S39}\label{maps:eq:E.S39}
\]

Here are the full dimensions for this small rational
calculation. Let $d_k$ be twice the degree-$k$ dimension
of $H^*(E)\otimes H^*(T^3)$ and let
$r_k=\dim(U_1^k+U_2^k)$:
\[
 \begin{array}{c|rrrrrrrrr}
 k&0&1&2&3&4&5&6&7&8\\ \hline
 d_k&2&8&20&44&70&68&36&8&0\\
 r_k&1&8&19&34&47&46&27&8&0
 \end{array}
 \tag{E.S40}\label{maps:eq:E.S40}
\]
Contraction in \eqref{maps:eq:E.S39} has rank one only for
$k=4,5$ and is zero in every other degree. All entries
are computed by exterior products of \eqref{maps:eq:E.S38} in
\texttt{maps\_shifted\_return.py}; its output retains an
explicit nonzero boundary witness in each of degrees
four and five. One can normalize the remaining sign in
$\phi_2$ by the ring structure: with
$r_1(1)=\theta+12\gamma$, the equality
$r_1(1)r_2(\delta)=r_3(a)$ forces the coefficient
$c_2=-2$, and then $c_3=2$, in the notation of \ref{maps:src:E.2.6}.
This is also the normalization used by that executable.
No rank calculation treats a rational cokernel as an
integral lattice.

\emph{The finite-cover cohomology needed for the full Betti numbers.}
For clarity, the full rational cohomology of $E_{12}$
can also be obtained without an uncomputed extension.
Let $\mathcal L_\lambda$ be the rank-one system on $E$
whose meridian acts by $\lambda$, with $\lambda^{12}=1$.
On a regular torus it is trivial. Its three base transports
in vertical degree $q$ are
\[
 \lambda^4\bigwedge^q A_1^{-t},\qquad
 \lambda^{-3}\bigwedge^q A_2^{-t},\qquad
                         \bigwedge^q M_0^{-t}.
 \tag{E.S41}\label{maps:eq:E.S41}
\]
The exponents retain the actual affine meridian character;
using only the finite linear matrices would give a different
cover.

At a finite cap, averaging on its finite root cover
identifies the cap cohomology with the invariant part of
\eqref{maps:eq:E.S41}. At the unchanged cusp the system is trivial,
because $g$ has character twelve and $u$ has character zero.
Its rational specialization is an isomorphism onto the
monodromy invariants in every degree. Here is an explicit
check of that last input. The cusp Betti numbers are
$(1,2,4,2,1)$ by its normalization calculation. In degree
two, the integral surjectivity of the specialization on
homology in \eqref{maps:eq:B.5.16} gives injectivity on cohomology;
the monodromy-invariant subspace has the same dimension
four. Its image therefore includes $uw$ and $\gamma\delta$.
In degree one the image is $\langle\gamma,u\rangle$;
their products with those two degree-two classes give the
two independent invariant classes in degree three and
the top class in degree four. Thus all the specialization
maps have the asserted ranks, with no additional
skyscraper summand. The paragraph following \eqref{maps:eq:B.5.16}
also proves these specialization assertions integrally
in all degrees by the actual honeycomb quotient map.

It follows that $R^q f_*\mathcal L_\lambda=j_*L_q$ over
the disc containing the three critical values. On this
disc $H^p(j_*L_q)=0$ for $p\geq2$. To check this directly,
retract the punctured disc to a free graph on three loops.
Restriction of a graph cocycle to the three local
coinvariant spaces is onto: its values on the three free
generators can be chosen independently. The exact
restriction sequence therefore kills $H^2(j_*L_q)$.
If $d_q=\binom4q$, $g_q$ is the common fixed dimension of
\eqref{maps:eq:E.S41}, and $k_{i,q}$ are their three individual fixed
dimensions, this same graph calculation gives
\[
 \dim H^0(j_*L_q)=g_q,\qquad
 \dim H^1(j_*L_q)=2d_q+g_q-\sum_i k_{i,q}.
 \tag{E.S42}\label{maps:eq:E.S42}
\]
The Leray sequence has only these two base degrees and
degenerates. Exact exterior-matrix ranks give the following
table; a row applies to each character of the indicated
order:
\[
 \begin{array}{c|rrrrrr}
 \operatorname{ord}(\lambda)&b_0&b_1&b_2&b_3&b_4&b_5\\ \hline
 1 &1&1&4&6&4&0\\
 2 &0&0&4&6&4&0\\
 3 &0&0&3&4&3&0\\
 4 &0&0&3&4&3&0\\
 6 &0&1&5&6&5&1\\
 12&0&1&4&4&4&1
 \end{array}
 \tag{E.S43}\label{maps:eq:E.S43}
\]
Here $b_j=\dim H^j(E;\mathcal L_\lambda)$ and $b_6=0$.
For the nontrivial characters all $g_q$ are zero; for the
trivial one they are all one. The program
\texttt{maps\_shifted\_cover.py} computes \eqref{maps:eq:E.S41}--\eqref{maps:eq:E.S43}
over the exact cyclotomic fields, represented by rational
companion matrices. It prints every local fixed dimension.
Summing the rows with multiplicities $1,1,2,2,2,4$ gives
\[
             (b_0,\ldots,b_6)(E_{12})=(1,7,46,56,46,6,0).
 \tag{E.S44}\label{maps:eq:E.S44}
\]
Both sides have Euler characteristic $24$, as required by
the degree-twelve covering of $E$.

\emph{The Betti and Jordan calculations on the actual fibre.}
The return on each of the two pieces of \eqref{maps:eq:E.S29}, up to
isotopy allowing its boundary to move, is its cyclic deck
transformation times the identity of $T^3$. This follows
also by moving the interpolation in \eqref{maps:eq:E.S36} entirely into
the other piece's collar. It acts trivially on the
cohomology of $D$. The generalized eigenvalue-one part of
the Mayer--Vietoris sequence therefore uses just the
invariant part $H^*(E)$ in \eqref{maps:eq:E.S40}. Its dimensions are
\[
       (1,0,1,19,45,45,19,1,0,1).
 \tag{E.S45}\label{maps:eq:E.S45}
\]
Every other eigenspace is the direct sum of the two piece
eigenspaces, because the boundary has none. In particular
those eigenspaces are semisimple. Their combined Poincare
polynomial is
\[
     2(1+t)^3(6t+42t^2+50t^3+42t^4+6t^5),
 \tag{E.S46}\label{maps:eq:E.S46}
\]
by \eqref{maps:eq:E.S43}--\eqref{maps:eq:E.S44}. Adding \eqref{maps:eq:E.S45} proves \eqref{maps:eq:E.S33}.

Finally the Wang sequence for $P=M_\psi$ and the rational
Betti numbers in \ref{maps:src:E.2.7} give the numbers of Jordan blocks
at eigenvalue one as
\[
          (1,0,1,19,44,44,19,1,0,1).
 \tag{E.S47}\label{maps:eq:E.S47}
\]
The only possible block sizes are one and two: the
twelfth power is \eqref{maps:eq:E.S31}, and \eqref{maps:eq:E.S37} squares to zero.
Subtracting \eqref{maps:eq:E.S47} from \eqref{maps:eq:E.S45} gives exactly one block
of size two in each of degrees four and five. This also
independently checks the direct variation ranks \eqref{maps:eq:E.S39}.
All other eigenvalues and blocks are already described
by \eqref{maps:eq:E.S43} and \eqref{maps:eq:E.S46}. None of these rational calculations
claims an integral torsion classification or replaces the
smooth $E_{12}$ by a manifold with the same invariants.
\end{proof}


